\question[6]
Liste ist unser selbstgebastelter Datentyp. Benutze nur die Methoden der Liste.
\begin{parts}
\part Erzeuge eine Liste a mit den Zahlen von 1 bis 100 in aufsteigender Reihenfolge.
\part Berechne die Summe der Elemente von a.
\part Verschiebe alle durch 3 teilbaren Zahlen aus a in eine neue Liste b. Die Reihenfolge der Zahlen soll
erhalten bleiben. Danach ist die Summe der Elemente von a gleich 3367 und die von b gleich 1683.
\end{parts}
\begin{solutionbox}{9.5cm}
\begin{lstlisting}
a = Liste()
for i in range(1, 101):
    a.insert(i)
    a.advance()

summe = 0
a.reset()
while not a.endpos():
    summe += a.elem()
    a.advance()

b = Liste()
a.reset()
while not a.endpos():
    x = a.elem()
    if x % 3 == 0:
        b.insert(x)
        b.advance()
        a.delete()
    else:
        a.advance()
\end{lstlisting}
\end{solutionbox}
